\documentclass[11pt]{article}
% Omit random trailer identifiers for reproducible, text-auditable PDF metadata.
\pdftrailerid{}
\usepackage[T1]{fontenc}
\usepackage[utf8]{inputenc}
\usepackage{lmodern}
\usepackage[margin=1in]{geometry}
\usepackage{amsmath,amssymb,amsthm,mathtools}
\usepackage{enumitem,booktabs,array,longtable}
\usepackage{xcolor,microtype}
\usepackage{etoolbox}
\usepackage{xurl}
\usepackage{tikz}
\usetikzlibrary{arrows.meta,positioning}
\usepackage[colorlinks=true,linkcolor=blue!45!black,citecolor=green!35!black,urlcolor=blue!55!black]{hyperref}
\usepackage[capitalize,nameinlink]{cleveref}
\hypersetup{pdftitle={Conditional Elimination and Circular Storage Codes},pdfauthor={Hans Brende and GPT-6 Astra (OpenAI)},pdfsubject={Weighted storage bounds for circular targets, October 2026},pdfkeywords={storage codes, polymatroids, circular cliques, conditional elimination}}
\numberwithin{equation}{section}
% ---------------------------------------------------------------------------
% Theorem environments
% ---------------------------------------------------------------------------
\theoremstyle{plain}
\newtheorem{theorem}{Theorem}[section]
\newtheorem{lemma}[theorem]{Lemma}
\newtheorem{corollary}[theorem]{Corollary}
\newtheorem{proposition}[theorem]{Proposition}
\newtheorem{conjecture}[theorem]{Conjecture}
\newtheorem{question}[theorem]{Question}
\newtheorem{problem}[theorem]{Problem}
\theoremstyle{definition}
\newtheorem{definition}[theorem]{Definition}
\newtheorem{remark}[theorem]{Remark}
\newtheorem{example}[theorem]{Example}

% The Lean proof is the proof; the text gives a faithful outline of it.
\newenvironment{outline}[1][Outline of the formal proof]{\begin{proof}[#1]}{\end{proof}}
% Results that are not formalized get complete proofs, labelled as such.
\newenvironment{informalproof}[1][Proof (not formalized)]{\begin{proof}[#1]}{\end{proof}}

% ---------------------------------------------------------------------------
% Lean references
%   \lean{name}            a declaration, relative to the namespace CircularMemo
%   \formal{file}{names}   the formal counterpart of a numbered result: a file relative to
%                          lean/CircularMemo/ and a comma-separated list of declarations
%   \partformal{file}{names}{gap}
%                          a numbered result whose formal counterpart covers all but the
%                          part described in {gap}; the gap is proved or cited in the text
%   \notformalized         a numbered result that has no formal counterpart
%   \cited{where}          a result of an earlier paper used without proof: where it is proved,
%                          and whether it is Lean-verified there
% ---------------------------------------------------------------------------
% Lean names are typeset like paths, with line breaks allowed after "_", "." and "/".
\DeclareUrlCommand\leanname{\def\UrlFont{\ttfamily\upshape}}
\newcommand{\lean}[1]{\leanname{#1}}
\newcommand{\leanfile}[1]{\leanname{#1}}
\newcommand{\leanlist}[1]{\begingroup\def\leansep{}%
  \renewcommand*{\do}[1]{\leansep\leanname{##1}\def\leansep{, }}\docsvlist{#1}\endgroup}
\newcommand{\formal}[2]{\par\noindent{\small\upshape\raggedright\textsf{Lean:} \leanlist{#2}\unskip
  \nobreak\hfil\penalty50\hskip1em\hbox{}\nobreak\hfill[\leanfile{#1}]\parfillskip=0pt\par}%
  \nobreak\smallskip\nobreak\noindent\ignorespaces}
\newcommand{\partformal}[3]{\par\noindent{\small\upshape\raggedright\textsf{Lean (partial):} \leanlist{#2}\unskip
  \nobreak\hfil\penalty50\hskip1em\hbox{}\nobreak\hfill[\leanfile{#1}]\parfillskip=0pt\par}%
  \nobreak\noindent{\small\upshape\textsf{Not formalized:} #3}\par\nobreak\smallskip\nobreak\noindent\ignorespaces}
\newcommand{\notformalized}{\par\noindent{\small\upshape\textsf{Not formalized;} a complete proof is provided.}\par\nobreak\smallskip\nobreak\noindent\ignorespaces}
\newcommand{\cited}[1]{\par\noindent{\small\upshape\textsf{Cited:} #1}\par\nobreak\smallskip\nobreak\noindent\ignorespaces}

% ---------------------------------------------------------------------------
% Notation
% ---------------------------------------------------------------------------
\newcommand{\N}{\mathbb{N}}
\newcommand{\Z}{\mathbb{Z}}
\newcommand{\Q}{\mathbb{Q}}
\newcommand{\R}{\mathbb{R}}
\newcommand{\C}{\mathbb{C}}
\newcommand{\F}{\mathbb{F}}
\newcommand{\T}{\mathsf{T}}
\newcommand{\clog}[1]{\lceil \log_2 #1 \rceil}
\newcommand{\abs}[1]{\lvert #1 \rvert}
\newcommand{\card}[1]{\lvert #1 \rvert}
\newcommand{\eps}{\varepsilon}
\newcommand{\prev}{\mathrm{prev}}
\newcommand{\cost}{\mathrm{cost}}
\newcommand{\gn}{\mathrm{gn}}
\newcommand{\Sh}{\mathrm{Sh}}
\newcommand{\cl}{\mathrm{cl}}
\newcommand{\Aut}{\mathrm{Aut}}
\newcommand{\minrk}{\mathrm{minrk}}
\newcommand{\MAIS}{\mathrm{MAIS}}
\newcommand{\Crux}{\mathrm{Crux}}
\newcommand{\TCP}{\mathrm{TCP}}
\DeclareMathOperator{\rank}{rank}
\DeclareMathOperator{\dist}{dist}
\DeclareMathOperator{\poly}{poly}

\DeclareMathOperator{\tr}{tr}
\DeclareMathOperator{\supp}{supp}
\newcommand{\norm}[1]{\lVert#1\rVert}

\setlist{itemsep=3pt,topsep=5pt}
\setlength{\emergencystretch}{2em}
\title{Conditional Elimination and Circular Storage Codes}
\author{Hans Brende \and GPT-6 Astra (OpenAI)}
\date{October 2026}

\begin{document}
\maketitle
\begin{abstract}
A conditional-elimination inequality gives the weighted storage bound
$2h(V)\le\sum_v h(\{v\})$ on every Andrasfai graph and every circular
clique of ratio below three. This strictly extends the range at most
$14/5$ proved in the accompanying fourth paper. The new family bound,
homomorphism transfer, and common-alphabet cardinality corollary are
verified in Lean, using that paper's existing definitions. An exact
conditional-mutual-information identity and the unequal-alphabet
corollary have complete written proofs. The result does not supply
these graph hypotheses for general native tape cones.
\end{abstract}
\begin{quote}\small
\textbf{Contributions.} Hans Brende directed the research. GPT-6 Astra
(OpenAI) developed the new argument, formalization, and text under his
direction. The memo reuses the fourth paper's framework and results,
by Hans Brende and Claude Opus 5.5 (Anthropic).
\end{quote}
\section{The strengthening and its scope}

The fourth paper proves that every storage-valid polymatroid on a
circular clique $K_{p/q}$ with $p/q\le14/5$ satisfies a weighted half
bound~\cite[Proposition ``Further good targets'']{BrendeCurved}.
This memo extends that range to $p/q<3$. The exact endpoint fails on
$K_3$. The existing paper and its Lean declarations remain unchanged;
the memo supplies a general proof and derives the earlier range from it.

A polymatroid here is a function $h:2^V\to\R$ on a finite set, with
$h(\varnothing)=0$, monotonicity, and submodularity. Write
\[
 h(A\mid S)=h(A\cup S)-h(S),\qquad
 I_h(A;B\mid S)=h(A\mid S)+h(B\mid S)-h(A\cup B\mid S).
\]
The mutual information is nonnegative for the disjoint sets used below.
Validity for recovery sets $N(v)$ means $h(v\mid N(v))=0$.
Diminishing returns preserves this equality after enlarging the
conditioning set. For a graph, $N(v)$ is its neighborhood.
No representability, integrality, or equality of singleton values is assumed.

\section{Conditional elimination}

\begin{theorem}[An exact partition identity]\label{thm:elimination}
\notformalized
Let $P_1,\ldots,P_m$ partition $V$, where $m\ge2$ and empty parts are
allowed. Enumerate $V$ as $v_0,\ldots,v_{n-1}$ and put
$S_t=\{v_0,\ldots,v_{t-1}\}$, $P_i^t=P_i\setminus S_t$.
Let $a(t)$ be the part of $v_t$. Suppose there is $b(t)\ne a(t)$ with
\[
 h(v_t\mid S_t\cup P_{b(t)}^t)=0.
\]
Then
\begin{equation}\label{eq:elimination}
 \sum_i h(P_i)-2h(V)
 =\sum_{t=0}^{n-1}\ \sum_{i\notin\{a(t),b(t)\}}
                       I_h(v_t;P_i^t\mid S_t)\ge0.
\end{equation}
In particular $2h(V)\le\sum_i h(P_i)\le\sum_{v\in V}h(\{v\})$.
\end{theorem}
\begin{informalproof}
Set $J_t=\sum_i h(P_i^t\mid S_t)-2h(V\setminus S_t\mid S_t)$.
Fix a step and abbreviate its vertex, own part, supplying part, and
conditional entropy by $v,a,b,u=h(v\mid S_t)$.
The own-part contribution decreases by $u$. Each other part decreases
by $I_h(v;P_i^t\mid S_t)$, and the whole remaining entropy decreases
by $u$. Thus
\[
 J_t-J_{t+1}=\sum_{i\ne a}I_h(v;P_i^t\mid S_t)-u.
\]
The supplying-part condition makes its summand exactly $u$.
Cancel it and telescope, using $J_n=0$ and
$J_0=\sum_i h(P_i)-2h(V)$. Nonnegative mutual information and
subadditivity within the parts give the two inequalities.
The empty ground set and zero conditional entropies cause no exception.
\end{informalproof}

More explicitly, the weighted slack is the right side of
\eqref{eq:elimination} plus
$\sum_i(\sum_{v\in P_i}h(\{v\})-h(P_i))$, another sum of nonnegative
mutual informations. If the supplying condition is dropped, the same
calculation subtracts $\sum_t h(v_t\mid S_t\cup P_{b(t)}^t)$ from
the right side. That is an exact defect identity, not an assertion
that a low-defect order always exists.

The following equivalent style of elimination is the short formal
interface used for the graph application. It is stated separately
because Lean verifies its inequality, not the preceding general
multi-part identity.

\begin{lemma}[A three-record criterion]\label{lem:three-part}
\formal{Elimination.lean}{three_part_elimination}
Let $A,B,C\subseteq V$, not necessarily disjoint. Suppose that for
every proper revealed set $S\subsetneq V$ there is $v\notin S$ for
which one of the following holds:
\[
 \begin{array}{ll}
 v\in A,&h(v\mid S\cup B)=0;\\
 v\in B,&h(v\mid S\cup C)=0;\\
 v\in C,&h(v\mid S\cup A)=0.
 \end{array}
\]
Then $2h(V)\le h(A)+h(B)+h(C)$.
\end{lemma}
\begin{outline}
Induct downwards on the number of unrevealed vertices to prove
$2h(V)+h(S)\le h(S\cup A)+h(S\cup B)+h(S\cup C)$.
Adding an eligible $v$ changes neither the own-record nor the
supplying-record entropy. Submodularity bounds the increase of the
third record by $h(S\cup\{v\})-h(S)$.
The induction therefore passes from $S\cup\{v\}$ to $S$.
At $S=V$ it is equality; at $S=\varnothing$ it is the conclusion.
\end{outline}

For a proper partition into three independent graph parts $A,B,C$,
orient edges cyclically $A\to B\to C\to A$. If this orientation is
acyclic, it satisfies the criterion: in each unrevealed induced graph
choose a source. Its unrevealed neighbors all belong to the next part;
its other neighbors are already in $S$. Validity and diminishing
returns give the displayed dependence. This passage from an arbitrary
acyclic orientation is a written argument. The specific order for the
Andrasfai application below is checked directly in Lean.

Acyclicity cannot be omitted from the three-record claim. On a six-cycle
colored $A,B,C,A,B,C$, the binary word
$(U,V,U+V,U,V,U+V)$ for independent fair bits is recoverable at every
vertex. Its joint entropy is two, but each of the three color records
has entropy one. This does not violate the weighted half bound for
the six-cycle; it distinguishes that bound from this particular
three-record inequality.

\section{All circular ratios below three}

The circular clique $K_{p/q}$ has vertices $0,\ldots,p-1$, with
adjacency given by $q\le(j-i)\bmod p\le p-q$.
For $k\ge2$ put $\operatorname{And}(k)=K_{(3k-1)/k}$.

\begin{theorem}[All Andrasfai graphs]\label{thm:andrasfai}
\formal{Circular.lean}{rateBound_andrasfai}
Every valid polymatroid on $\operatorname{And}(k)$, $k\ge2$, satisfies
$2h(V)\le\sum_{v\in V}h(\{v\})$.
\end{theorem}
\begin{outline}
Use independent intervals
$A_i=i$, $B_j=k+j$ for $0\le i,j<k$, and
$C_\ell=2k+\ell$ for $0\le\ell<k-1$.
Directly from circular adjacency,
\[
 A_i\sim B_j\iff i\le j,\quad
 B_j\sim C_\ell\iff j\le\ell,\quad
 C_\ell\sim A_i\iff\ell<i.
\]
The order $A_0,B_0,C_0,A_1,B_1,C_1,\ldots,A_{k-1},B_{k-1}$
places every cyclically oriented edge forward. Lean checks these
inequalities using the integer order keys $3i,3j+1,3\ell+2$.
For any unrevealed set take its least key; all later neighbors lie
in the next part. \Cref{lem:three-part} gives
$2h(V)\le h(A)+h(B)+h(C)$, and subadditivity over the partition
gives the weighted assertion.
\end{outline}

\begin{corollary}[Circular targets and homomorphisms]\label{cor:circular}
\formal{Circular.lean}{rateBound_circularClique_lt_three,rateBound_of_circularHom,rateBound_circularClique_oldRange}
For integers $q>0$ and $0\le p<3q$, every valid polymatroid on
$K_{p/q}$ satisfies the weighted half bound. The same holds on every
finite simple graph with a homomorphism to such a clique.
The previous range $5p\le14q$, including its degenerate $p=q=0$
case under the Lean empty-graph convention, is a corollary.
\end{corollary}
\begin{outline}
Choose $k=q+1\ge2$. Since $p\le3q-1$,
$pk\le(3k-1)q$. The already formalized Bondy--Hell map
$i\mapsto\lfloor(3k-1)i/p\rfloor$ gives
$K_{p/q}\to\operatorname{And}(k)$~\cite{BondyHell1990}.
For $p=0$ this is the unique empty-domain map; no division is evaluated
at a vertex. Apply \cref{thm:andrasfai} and the fourth paper's
homomorphism pullback. That pullback sends $h$ to
$Z\mapsto h(\phi^{-1}(Z))$ and uses subadditivity in each fibre.
If $q>0$ and $5p\le14q$, then $p<3q$; if $q=0$, the older range
forces $p=0$ and the conclusion is the empty-graph bound.
No earlier numerical circular-clique certificate is used in this proof.
\end{outline}

\begin{corollary}[Arbitrary common alphabets]\label{cor:codes}
\formal{Codes.lean}{card_sq_le_of_circularHom}
Let $G$ be a finite simple graph mapping to $K_{p/q}$ with $q>0$
and $p<3q$. For any nonempty finite alphabet $\mathcal A$, every
storage code $\mathcal C\subseteq\mathcal A^{V(G)}$ satisfies
$|\mathcal C|^2\le|\mathcal A|^{|V(G)|}$.
No linearity assumption is imposed.
\end{corollary}
\begin{outline}
Supply \cref{cor:circular} to the fourth paper's existing
cardinality theorem. That theorem applies the
weighted bound to the entropy of a uniform codeword and converts the
logarithmic inequality to the stated integer one. It handles empty
codes and one-symbol alphabets as well. No entropy library is duplicated.
\end{outline}

\begin{corollary}[Unequal alphabets]\label{cor:unequal}
\notformalized
On the same graphs, every storage code
$\mathcal C\subseteq\prod_v\mathcal A_v$ over nonempty finite
alphabets satisfies $|\mathcal C|^2\le\prod_v|\mathcal A_v|$.
\end{corollary}
\begin{informalproof}
For a nonempty code, take its uniform law. Its entropy is a valid
polymatroid and has $h(V)=\log|\mathcal C|$ and
$h(\{v\})\le\log|\mathcal A_v|$. Apply \cref{cor:circular} and
exponentiate. Empty codes satisfy the inequality directly.
\end{informalproof}

The endpoint is sharp: $K_{3/1}=K_3$ supports the binary parity code
$(U,V,U+V)$ of size four, whose square exceeds $2^3$.
The half constant is also attained on each Andrasfai graph with a
suitable alphabet: give each edge an independent $a$-ary symbol and
store its incident tuple at each vertex. The graph is $k$-regular,
so the coordinate alphabet has size $a^k$, and
$2|E|=k|V|$ gives rate exactly one half. This is the standard edge-symbol
construction, not a new encoding method.

\paragraph{Contribution to the program.}
The memo widens an all-alphabet impossibility tool and replaces isolated
circular-target certificates by one explicit elimination mechanism.
A prospective construction on these graphs cannot exceed half rate.
An application to native transposition still needs a proved graph
homomorphism or elimination interface for its actual recovery relation;
neither is furnished here. The general tape question remains open.
The argument's correctness and formal verification are separate from
external priority, which has not been established for this strengthening.

\appendix
\section{Index of Lean declarations}\label{sec:leanindex}

Declarations are relative to \texttt{CircularMemo}; their axiom checks are generated together with this index (\cref{sec:verification}).

\paragraph{\texttt{CircularMemo.Circular}}
{\small\raggedright \texttt{rateBound\_andrasfai}, \texttt{rateBound\_circularClique\_lt\_three}, \texttt{rateBound\_circularClique\_oldRange}, \texttt{rateBound\_of\_circularHom}.\par}

\paragraph{\texttt{CircularMemo.Codes}}
{\small\raggedright \texttt{card\_sq\_le\_of\_circularHom}.\par}

\paragraph{\texttt{CircularMemo.Elimination}}
{\small\raggedright \texttt{three\_part\_elimination}.\par}

\section{Verification and reuse}\label{sec:verification}

The authoritative proofs of the results marked \textsf{Lean} are in
the memo's \texttt{lean/CircularMemo/} library. The English outlines
describe those proofs for convenience. The general partition identity,
the general acyclic-orientation argument, unequal alphabets, and the
sharpness examples have written proofs; they are not claimed to be
formalized by the three-part inequality alone.

The library requires the fourth paper by the relative path
\texttt{../../lean}; that paper in turn requires the earlier shared
libraries. Clone the whole repository, preserving sibling directories.
Lean and Mathlib remain pinned to v4.34.1. No new axiom, proof placeholder,
or kernel-bypassing decision procedure is used. The permitted axiom set
is \texttt{propext}, \texttt{Classical.choice}, and \texttt{Quot.sound}.

From the memo directory, run \texttt{python3 tools/lean\_index.py}.
It generates this declaration index and \texttt{lean/PaperAxioms.lean}
from the same formal-status annotations. In \texttt{lean/}, run
\texttt{lake build} and \texttt{lake env lean PaperAxioms.lean}.
On a fresh clone first fetch the pinned dependency cache with
\texttt{lake exe cache get}. In \texttt{paper/}, run
\texttt{latexmk -pdf main.tex}.

The main paper's original declarations and certificate remain intact.
The memo's old-range corollary derives that range directly from the
stronger theorem, including the empty case.
Removing the old certificate from its original module is deferred until
a separate dependency and publication review.

\bibliographystyle{plainurl}
\bibliography{references}
\end{document}
